\begin{table}[!htbp]
\captionsetup{labelsep=period}
\caption{Definitions used in the review. Evidence must cover both methods and the stated main comparison. These are operational criteria for this review.}
\label{tab:peft-review-criteria}
\centering
\begin{tabular}{@{}p{.25\linewidth}p{.71\linewidth}@{}}
\toprule
Control & Required evidence \\
\midrule
Training runs and uncertainty & At least two independent training runs per method with a defined SD, SE or CI attributable to training variation. \\
LR selection & At least three learning rates evaluated per method, with separate selection using a stated validation criterion. \\
LR candidates & Numerical candidate values or a search range stated for both methods. \\
Comparable conditions & Shared checkpoint and data, comparable exposure, and adapted modules and trainable budgets matched or addressed through an appropriate controlled comparison. \\
\bottomrule
\end{tabular}
\end{table}

\begin{table}[!htbp]
\captionsetup{labelsep=period}
\caption{Paper-level coverage of experimental controls. The first three columns partition the 64 eligible papers using the at least one comparison rule. The final column requires established coverage of every applicable main comparison. Unresolved coverage prevents a universal positive.}
\label{tab:peft-review-coverage}
\centering
\begin{tabular*}{\linewidth}{@{\extracolsep{\fill}}lrrrr@{}}
\toprule
Control & Documented & Explicit departure & Unclear & \shortstack{All main\\comparisons} \\
\midrule
Training runs and uncertainty & 9 & 5 & 50 & 2 \\
LR selection & 7 & 4 & 53 & 1 \\
LR candidates & 13 & 0 & 51 & 4 \\
Comparable conditions & 25 & 10 & 29 & 1 \\
\bottomrule
\end{tabular*}
\end{table}

\begin{table}[!htbp]
\captionsetup{labelsep=period}
\caption{Documented controls within each domain. Entries give qualifying papers over applicable papers. A paper can occur in several domains, so rows cannot be summed. The decoder subset restricts comparisons to standard decoder adaptation against vanilla LoRA, with the exclusions described by Lee et al. Its definitions and sampling frame still differ from theirs.}
\label{tab:peft-review-domains}
\centering
\begin{tabular*}{\linewidth}{@{\extracolsep{\fill}}lrrrr@{}}
\toprule
Domain & \shortstack{Training\\uncertainty} & LR selection & LR candidates & Conditions \\
\midrule
Decoder language & 4/51 & 4/51 & 7/51 & 18/51 \\
Other language & 5/40 & 3/40 & 6/40 & 7/40 \\
Vision and multimodal & 3/20 & 1/20 & 3/20 & 2/20 \\
Diffusion & 0/11 & 0/11 & 1/11 & 3/11 \\
Other & 0/1 & 0/1 & 0/1 & 0/1 \\
\midrule
Decoder subset & 3/40 & 2/40 & 4/40 & 12/40 \\
\bottomrule
\end{tabular*}
\end{table}
